\section{Presupuesto de margen y contratos de retirada (INT-31)}
\label{sec:sd4-integracion-seguridad31}
Se fija para el cribado de la cadena una reserva propia de generación del 25\,\% y una reducción propia del caudal útil del 10\,\% respecto de su indicación. Son condiciones de diseño que deben cubrirse mediante selección, calibración y dispersión; no son tolerancias OEM, factores normativos ni mediciones de emisión. No sustituyen la clasificación de área, la ventilación distribuida ni el expediente de seguridad. Se conserva la relación conservadora de GAS-28, $k=7,5978/60=0,12663$ m$^3$/h por amperio agregado de cadenas, y volumen de proyecto 36 m$^3$.

Con indicación de 900 m$^3$/h y corriente real agregada de 20 A, el cribado resulta $Q_u=810$ m$^3$/h, $q_d=1,25(0,12663)(20)=3,16575$ m$^3$/h y $100q_d/Q_u=0,390833\,\%$. La corriente real máxima sin margen adicional para no alcanzar 0,4\,\% es $0,004(810)/(1,25\cdot0,12663)=20,4691$ A. El mínimo útil para 20 A es 791,4375 m$^3$/h, equivalente a indicación de 879,375 m$^3$/h bajo aquella reducción. Por ello 900/20 deja un margen pequeño: no se convierte en ajuste habilitado. Un error superior de corriente del 5\,\%, si fuese el presupuesto seleccionado, llevaría 20 A indicados a 21 A reales y 0,410375\,\%, por encima de prealarma. Con ese error la indicación límite sería 19,4944 A antes de otras incertidumbres. No se atribuye ese 5\,\% al cargador ni se impone un ajuste que su resolución no permita. El expediente debe conciliar límite real autorizado, dispersión entre cadenas, medida y caudal útil; la selección de 20 A nominales sin esos datos no cierra RT-06.07/11.

\textbf{Variables y separación funcional.} El permiso de recarga requiere por cada banco medida válida de corriente agregada de todas sus cadenas, caudal útil de extracción válido, gas válido, condiciones térmicas y circuito final disponible. La medida de impulsión DPT no se reutiliza como prueba de extracción; la medida SEL de grupo no se transforma en corriente DC. Cada variable se registra con unidad, rango, incertidumbre, límite, validez y edad de muestra. Una lectura fuera de rango, congelada o perdida retira el permiso; una señal de comunicación recibida no acredita caudal ni corriente. La realización física debe seleccionar entradas, diagnóstico, alimentación y contactos antes de habilitar. Se conserva el circuito local de retirada independiente de nube y WAN.

\textbf{Estados de proyecto.} Al energizar o tras una parada, el estado inicial es recarga inhibida. Se comprueban gas y extracción antes de armar; el rearme manual reconoce una causa resuelta, sin puentear señal inválida ni disparo enclavado. Pérdida de caudal, instrumento, alimentación o exceso retira permiso y enclava investigación. Retirar AC del rectificador puede llevar la carga a batería: el estado se registra y debe conservar ambiente y reserva. La recuperación de caudal, tensión o comunicación no repone recarga automáticamente. El permiso de bancada del generador conserva separación de SR-R: retirar carga resistiva no acredita que recarga y servicio caben en el grupo. El cierre contra incendio se coordina con la generación residual, sin considerar la retirada AC como emisión cero.

\textbf{Purga y sellado.} Desde 0,8\,\% hasta 0,2\,\%, con 810 m$^3$/h útiles y emisión exactamente cero, el modelo ideal exige $t=(36/810)\ln(0,008/0,002)=221,807$ s. Es una cota ideal de cribado, no temporizador de habilitación: emisión residual, estratificación y volúmenes muertos impiden certificar purga por tiempo. Con la generación de 20 A y reserva anterior, el equilibrio 0,390833\,\% ni siquiera permite llegar a 0,2\,\% por ventilación continua. Debe comprobarse el estado real tras aislamiento antes del rearme. Si el sellado inicia a 0,2\,\%, permanecer por debajo de 1,0\,\% durante diez minutos sin extracción requiere $q_{res}\leq(0,010-0,002)(36)/(10/60)=1,728$ m$^3$/h bajo mezcla perfecta. Ninguna ficha de aspirante ni lectura inicial demuestra ese residual. La retención de agente y el enclavamiento por H$_2$ quedan vinculados a esta verificación, sin autorizar sellado o descarga de banco mediante un retardo calculado.

La recepción prevista registra cada fallo de señal, retirada física, corriente de cadenas, caudal efectivo y evolución de gas, junto con reserva UPS y temperatura. Su informe distingue ensayo de un componente de prueba integral de la cadena. La secuencia de estados constituye especificación funcional de proyecto; no se presenta como programa cargado, ensayo realizado ni nivel de seguridad certificado. Continúan pendientes la selección de la cadena de medida/maniobra y las validaciones OEM y de conjunto.

